\subsection{The same quantum detector with reconstructed action-time uncertainty}
\label{sec:source-scalar-clock-quantum}

The original-vacuum probability comparison extends to the elapsed-time
intervals obtained from Proposition~\ref{prop:source-scalar-clock-stability}
on the same authenticated scalar records. Fix the
positive $64$-mode action, mass metric, original vacuum, coherent momentum
preparation $v$, and detector $g$ of
Proposition~\ref{prop:source-scalar-quantum-probability}.
The quantum state and operator are kept fixed throughout this comparison.
Uncertain record-based timing is the additional input.
The extension covers reference probabilities inside each time interval;
it does not define a continuous trajectory for the split evolution.

\begin{proposition}[Probability comparison on a recovered time interval]
\label{prop:source-scalar-clock-probability}
Let $p_j$ be the intervention probability of
$E_g=(I+\sin\Phi(g))/2$ after $j$ executed symmetric gates, and let $p(t)$
be its continuous finite-action probability on the same preparation.
Suppose $|p_j-p(t_j)|\le\epsilon_j$, where $t_j=j\tau$ is the nominal
action time. For any closed interval $I_j=[a_j,b_j]$, set
\[
 L_g=\tfrac12\|g\|_M\|v\|_M,\qquad
 \delta_j=\max\{|a_j-t_j|,|b_j-t_j|\}.
\]
Then, simultaneously for every $t\in I_j$,
\begin{equation}
 |p_j-p(t)|\le \epsilon_j+L_g\delta_j.
 \label{eq:source-scalar-clock-probability-bound}
\end{equation}
Both baseline probabilities are $1/2$, so this also bounds the difference
of the corresponding intervention responses. If a certified signed split
response has magnitude at least $s_j>\epsilon_j+L_g\delta_j$, the
continuous response has that same sign throughout $I_j$, with magnitude
at least $s_j-\epsilon_j-L_g\delta_j$.
\end{proposition}
\begin{proof}
The original vacuum is stationary under the continuous Hamiltonian.
Its detector variance $V_0$ is constant and the coherent field mean is
\[
 m(t)=\langle g,\mathsf A^{-1/2}\sin(t\sqrt{\mathsf A})v\rangle_M,
 \qquad
 m'(t)=\langle g,\cos(t\sqrt{\mathsf A})v\rangle_M.
\]
Self-adjoint spectral calculus gives
$\|\cos(t\sqrt{\mathsf A})\|_M\le1$ for every real $t$.
The finite-dimensional mean is differentiable, and
$p(t)=(1+e^{-V_0/2}\sin m(t))/2$ therefore obeys
$|p'(t)|\le e^{-V_0/2}\|g\|_M\|v\|_M/2\le L_g$.
The mean-value bound and triangle inequality prove
\eqref{eq:source-scalar-clock-probability-bound}.
The split probability belongs only to its executed layer; no split
evolution at unexecuted times is introduced by this argument.
\end{proof}

\paragraph{Exact comparison on all retained intervals.}
The two execution schedules have identical reconstructed field layers
and elapsed-time enclosures. The conditional clock budget is
$10^{-8}$ per configuration in the $M$ norm, with a common hidden
history stationary for the same supplied trapezoidal action and one
positive duration per complete configuration interval.
Exact source moments give $L_g\le0.092011181733$.
Outward rational arithmetic composes all $21$ original probability rows
with all $21$ positive elapsed-time intervals, retaining every unresolved
row. The additional timing error is at most $0.000000449304$.
The same $15$ responses are resolved uniformly over their
respective nonzero intervals.

For the sixteenth completed update, the recovered interval is
\[
 I_{16}=[1.022198489136,\,1.022206518861].
\]
The split probability remains in
$[0.490342191369,0.490493790733]$, and for every $t\in I_{16}$,
\[
 |p_{16}-p(t)|\le0.000811258283,\qquad
 p(t)-\tfrac12\le-0.008694950984.
\]
These are interval-wide predictions of the same finite quantum action.
The source, state, covariance and detector are unchanged.

The record-error budget conditions the clock inverse problem; it is not
an error bound for a perturbed quantum preparation, field trajectory,
Hamiltonian or detector. Existence of one common hidden stationary history
is a hypothesis and does not follow from intersecting duration intervals.
The graph reconstruction supplies complete configurations for this finite
record, not a regional quantum time-slice theorem. The action normalization,
variable-duration quadrature, orientation, vacuum and canonical
quantization remain declared. No count-volume identification, physical
seconds, quantum clock measurement, observed outcomes or spatial-continuum
error is inferred from the timing comparison.
